\chapter{Foundations}
\label{c:chapter1}

This chapter introduces the formal and conceptual background required for the remainder of the thesis. We begin with the mathematical foundations of network science and distinguish static from temporal contact networks, emphasising the causal structure that temporal ordering imposes on spreading processes. We then review the family of compartmental epidemic models, with particular attention to the stochastic SIR model that underpins all experiments in this work. Classical source detection algorithms are surveyed next, establishing the non-machine-learning baselines against which the proposed methods are later evaluated. The chapter concludes with a self-contained introduction to graph neural networks, covering the general message-passing paradigm, its principal static instantiations, and the challenges that arise when extending these methods to temporal network data.

\section{Network Science and Graph Theory}
\label{st:section1}

A graph $G = (\mathcal{V}, \mathcal{E})$ consists of a finite set of nodes $\mathcal{V}$ with $|\mathcal{V}| = N$ and a set of edges $\mathcal{E} \subseteq \mathcal{V} \times \mathcal{V}$. In an undirected graph, each edge $\{u, v\} \in \mathcal{E}$ represents a symmetric relationship; in a directed graph, edges are ordered pairs $(u, v)$ indicating a relationship from $u$ to $v$. The topology of a graph is compactly encoded by its adjacency matrix $\mathbf{A} \in \{0,1\}^{N \times N}$, where $A_{uv} = 1$ if and only if $(u, v) \in \mathcal{E}$. For undirected graphs, $\mathbf{A}$ is symmetric. The degree of a node $v$, denoted $d_v$, counts the number of edges incident to it; for directed graphs, one distinguishes the in-degree $d_v^{\mathrm{in}}$ and out-degree $d_v^{\mathrm{out}}$. The neighbourhood of $v$ is the set $\mathcal{N}(v) = \{u \in \mathcal{V} : (u,v) \in \mathcal{E}\}$ of nodes adjacent to $v$. Nodes may carry feature vectors $\mathbf{x}_v \in \mathbb{R}^d$, and edges may carry weights or attributes; these attributed graphs form the primary input to the GNN models studied in this thesis.

\subsection{Static vs. Temporal Networks}

A static network treats the edge set $\mathcal{E}$ as fixed over time. This is an appropriate abstraction for many structural or relational datasets, but it is too coarse for systems in which interactions are transient and sequentially ordered. A temporal network, by contrast, is a triple $\mathcal{G} = (\mathcal{V}, \mathcal{E}^t, \mathcal{T})$, where $\mathcal{T} = \{t_1, t_2, \ldots, t_M\}$ is an ordered set of event times and each edge $(u, v, t) \in \mathcal{E}^t$ records a contact between nodes $u$ and $v$ occurring at time $t$. Temporal networks may equivalently be represented as a sequence of graph snapshots $\{G_0, G_1, \ldots, G_T\}$, where $G_t = (\mathcal{V}, E_t)$ and $E_t$ contains all contacts active during the time interval $[t, t+1)$.

The aggregated static graph $G_a = (\mathcal{V}, E_a)$ is obtained by taking the union of all snapshot edge sets, $E_a = \bigcup_{t=0}^{T} E_t$. While $G_a$ records which pairs of nodes were ever in contact, it discards the temporal ordering entirely. This loss of information is consequential for epidemic spreading: a contact $(A, B, t=1)$ followed by $(B, C, t=2)$ permits the transmission chain $A \to B \to C$, whereas the reversed sequence $(B, C, t=1)$ followed by $(A, B, t=2)$ does not, because $B$ cannot have been infected by $A$ before the contact between them occurred. The static projection $G_a$ treats both orderings identically. Real-world temporal contact networks arise in many domains, including face-to-face proximity measured by RFID or Bluetooth sensors, digitised livestock movement records, and timestamped email or messaging exchanges.

\subsection{Causality and Causal Walks in Dynamic Graphs}

The causal structure of a temporal network is captured by the notion of a time-respecting path, also called a causal walk. Formally, a causal walk from node $v_0$ to node $v_k$ is a sequence of contacts $((v_0, v_1, t_1), (v_1, v_2, t_2), \ldots, (v_{k-1}, v_k, t_k))$ such that all contacts belong to $\mathcal{E}^t$ and the event times are strictly increasing, $t_1 < t_2 < \cdots < t_k$. The existence of a causal walk from $v_0$ to $v_k$ is the necessary condition for a pathogen originating at $v_0$ to reach $v_k$ through sequential transmission.

The reachable set $\mathcal{R}(v, T)$ from a source node $v$ by time $T$ is defined as the set of all nodes $w$ for which a causal walk from $v$ to $w$ exists with $t_k \leq T$. This set can be strictly smaller than the reachable set on the static projection $G_a$, particularly in networks with bursty or structured contact patterns. For source detection, the causal path structure is the fundamental quantity of interest: the set of candidate sources for an observed infected node $i$ at time $T$ is precisely the set of nodes from which a causal walk leading to $i$ exists before $T$. Methods that operate on $G_a$ rather than $\mathcal{G}$ implicitly treat both causally valid and invalid transmission chains as equally plausible, introducing false candidates into the inference problem. Exploiting causal walk structure is therefore not merely a modelling refinement but a logical necessity for exact source inference on temporal networks.

\section{Epidemic Spreading Models}
\label{st:section2}

Compartmental models are the dominant framework for mathematical epidemiology. They partition the population into a small number of discrete states and specify probabilistic rules governing transitions between states. The simplest models are fully determined by two parameters: a transmission probability $\beta$ governing the rate at which the pathogen spreads from infected to susceptible individuals, and a recovery probability $\mu$ governing the rate at which infected individuals recover. We describe the principal compartmental models in order of increasing complexity before specialising to the SIR model used throughout this thesis.

\subsection{The SI, SIR, SIS, and SIRS Models}

The SI model is the simplest epidemic model, containing only two compartments: susceptible (S) and infectious (I). An infected node transmits to each susceptible neighbour with probability $\beta$ per contact; there is no recovery. The entire connected component of the source eventually becomes infected, making SI appropriate only for processes where immunity or intervention is absent.

The SIS model extends SI by allowing recovery: an infected node transitions back to the susceptible state with probability $\mu$ per time step, reflecting the assumption that infection confers no lasting immunity. This produces an endemic equilibrium in which the pathogen persists within the population, and it is the standard model for recurrent illnesses such as the common cold. The SIR model introduces a third compartment, recovered (R): an infected node transitions to R with probability $\mu$, and recovered nodes are permanently immune and cannot be reinfected. The SIR model is appropriate for acute infectious diseases such as influenza or measles and is the primary model used in this thesis. The SIRS model allows immunity to wane over time, so that recovered nodes return to the susceptible pool with probability $\omega$ per time step, enabling recurrent epidemic waves. For the purposes of this work, the irreversibility of the recovered state in the SIR model is a key property: the final observed snapshot of SIR states constitutes a permanent record of which nodes were ever infected, simplifying the inference problem.

\subsection{The SIR Model and Stochastic Transitions}

We define the stochastic SIR process on a temporal contact network $\mathcal{G}$. Each node $v \in \mathcal{V}$ occupies a state $X_v(t) \in \{S, I, R\}$ at each discrete time step $t$. The process is initialised by setting $X_s(0) = I$ for the source node $s$ and $X_v(0) = S$ for all $v \neq s$. Two types of stochastic transition are applied at each time step. First, for each contact $(i, j, t) \in \mathcal{E}^t$ where $X_i(t) = I$ and $X_j(t) = S$, the susceptible node $j$ becomes infected with probability $\beta$:
\[
    \mathbb{P}\!\left(X_j(t+1) = I \mid X_j(t) = S,\ X_i(t) = I,\ (i,j,t) \in \mathcal{E}^t\right) = \beta.
\]
Second, each infected node $i$ with $X_i(t) = I$ independently recovers with probability $\mu$:
\[
    \mathbb{P}\!\left(X_i(t+1) = R \mid X_i(t) = I\right) = \mu.
\]
Recovered nodes remain in state R for all subsequent time steps. The process terminates when no infected nodes remain.

In this thesis, $\beta$ and $\mu$ are fixed hyperparameters specified in experiment configuration files. Different values of $(\beta, \mu)$ produce outbreaks of different expected size and duration: high $\beta$ relative to $\mu$ yields large outbreaks where many nodes are infected, while low $\beta$ or high $\mu$ produces small, rapidly extinguished outbreaks. Training data is generated by running the SIR process repeatedly from each candidate source, with independent random draws for all transition events at each repetition. This Monte Carlo procedure, implemented using an adapted version of Holme's fast C-based temporal SIR simulation engine~\cite{Holme2020}, produces large ensembles of simulated epidemic outcomes that form the supervised training and evaluation datasets for all models in this thesis.

\section{Classical Source Detection Algorithms}
\label{st:section3}

Before the advent of machine learning approaches to source detection, the field relied on two complementary families of methods: analytical estimators derived from probabilistic first principles, and simulation-based estimators that score candidate sources by comparing observed and simulated outbreak snapshots. We summarise the most relevant representatives of each family, as they constitute the classical baselines evaluated in Chapter~6.

\subsection{Rumor Centrality and Jordan Center}

Shah and Zaman~\cite{Shah2011} established the theoretical foundation of algorithmic source detection by showing that the source of a spreading process on a regular tree can be identified as a maximum likelihood estimator under the SI model. They defined the rumour centrality $R(v)$ of a node $v$ as the number of orderings in which the observed infection pattern can arise if $v$ is the source, and proved that the node maximising $R(v)$---the rumour centre---is a strongly consistent estimator of the true source as the network size grows. The Jordan centre extends this idea to general network topologies and the SIR model: it is defined as the node that minimises the maximum shortest-path distance to any infected or recovered node, $v^* = \arg\min_{v \in \mathcal{V}_{I \cup R}} \max_{u \in \mathcal{V}_{I \cup R}} d(v, u)$, where $d(\cdot, \cdot)$ denotes the shortest path on the induced infected subgraph. Both estimators are computationally inexpensive and require no knowledge of the spreading parameters $\beta$ or $\mu$. Their principal limitation is a reliance on tree-structure approximations and shortest-path distances on the static projection, which means they cannot exploit the temporal ordering of contacts or account for cycles in the contact network.

\subsection{Simulation-based Methods: Soft Margin Estimator (SME)}

The Soft Margin Estimator (SME) of Antulov-Fantulin et al.~\cite{AntulovFantulin2015} takes a Monte Carlo approach to source scoring. For each candidate source $v$, $M$ independent SIR simulations are run from $v$ on the contact network, producing a collection of simulated final snapshots $\{\tilde{r}_{v,m}\}_{m=1}^M$. The SME score for source $v$ is then the mean similarity between the observed snapshot $r^*$ and each simulated snapshot, measured by a soft-margin kernel:
\[
    P_{\mathrm{SME}}(v) \propto \frac{1}{M} \sum_{m=1}^{M} \exp\!\left(-\left[\varphi(r^*, \tilde{r}_{v,m}) - 1\right]^2\right),
\]
where $\varphi(r^*, \tilde{r}_{v,m})$ is the Jaccard similarity between the observed and simulated infected-plus-recovered sets. The SME is model-agnostic in the sense that any stochastic spreading process can be substituted in place of SIR, and it naturally handles partial observations by comparing only those nodes for which state information is available. Its principal weakness is computational cost: scoring all $N$ candidate sources requires $N \times M$ full epidemic simulations, which is prohibitive for large networks and typically demands $M \geq 1000$ simulations per candidate to achieve stable score estimates. As noted by Sterchi et al.~\cite{Sterchi2025}, SME is frequently used as a gold-standard classical baseline in benchmarks of GNN-based methods precisely because it is the closest classical analogue to the simulation-based training signal used by supervised neural approaches.

\section{Graph Neural Networks (GNNs)}
\label{st:section4}

Graph neural networks are a family of architectures that generalise deep learning to graph-structured data by defining differentiable operations that respect the topology of the input graph. Unlike convolutional neural networks, which exploit the regular grid structure of image data, GNNs must be permutation-equivariant with respect to node ordering and must aggregate information across irregular, variable-degree neighbourhoods. The message-passing framework provides the canonical abstraction for this class of models.

\subsection{Message Passing Paradigm}

The general message-passing neural network (MPNN) framework, formalised by Gilmer et al.~\cite{Gilmer2017}, describes GNN computation in terms of two alternating operations applied at each layer $l = 1, \ldots, L$. In the aggregation step, each node $v$ collects messages from its neighbours:
\[
    \mathbf{m}_v^{(l)} = \mathrm{AGG}\!\left(\left\{M_l\!\left(\mathbf{h}_v^{(l-1)},\, \mathbf{h}_u^{(l-1)},\, \mathbf{e}_{uv}\right) : u \in \mathcal{N}(v)\right\}\right),
\]
where $M_l$ is a learnable message function, $\mathbf{e}_{uv}$ is an optional edge feature, and $\mathrm{AGG}$ is a permutation-invariant aggregator such as sum, mean, or max. In the update step, the node's own embedding is combined with the aggregated message:
\[
    \mathbf{h}_v^{(l)} = U_l\!\left(\mathbf{h}_v^{(l-1)},\, \mathbf{m}_v^{(l)}\right).
\]
Node embeddings are initialised from input features, $\mathbf{h}_v^{(0)} = \mathbf{x}_v$. After $L$ layers of message passing, the embedding $\mathbf{h}_v^{(L)}$ encodes a summary of the $L$-hop neighbourhood of $v$. Different architectural choices for $M_l$, $\mathrm{AGG}$, and $U_l$ yield the principal GNN variants. The Graph Convolutional Network (GCN) of Kipf and Welling~\cite{Kipf2017} uses a symmetric degree-normalised aggregation without a distinct message function:
\[
    \mathbf{h}_v^{(l)} = \mathrm{ReLU}\!\left(\mathbf{W}^{(l)} \sum_{u \in \mathcal{N}(v) \cup \{v\}} \frac{\mathbf{h}_u^{(l-1)}}{\sqrt{\tilde{d}_v\,\tilde{d}_u}}\right),
\]
where $\tilde{d}_v = d_v + 1$ accounts for the added self-loop. GraphSAGE~\cite{Hamilton2017} concatenates the node's own embedding with the mean of its neighbours' embeddings before applying a linear transformation, making it well-suited for inductive settings. The Graph Attention Network (GAT) of Veli\v{c}kovi\'{c} et al.~\cite{Velickovic2018} replaces fixed normalisation with learned attention coefficients, allowing the model to weight neighbour contributions dynamically based on feature similarity. All three variants are evaluated as static baselines in this thesis.

\subsection{GNNs for Static Source Detection (GCN, GraphSAGE, GAT)}

The source detection problem on a static network can be formulated as a node classification task: given an attributed graph $(G_a, \mathbf{X})$ where $\mathbf{X} \in \{0,1\}^{N \times 3}$ encodes the one-hot SIR state of each node at the time of observation, predict which node is the epidemic source. GNNs are particularly well-suited to this task because source nodes tend to lie near the topological centre of the infected subgraph, and iterative neighbourhood aggregation naturally propagates the local density of infected nodes inward toward the source. This intuition is supported empirically by the benchmark of Sterchi et al.~\cite{Sterchi2025}, who show that GCN, GraphSAGE, and GAT all substantially outperform classical methods including the Jordan centre and the SME across a range of network topologies. Training is performed in a supervised manner on a large collection of simulated epidemic outcomes: each training sample is a triple consisting of the contact graph, the node-state feature matrix, and the identity of the true source node. The loss is the cross-entropy between the softmax output distribution over nodes and the one-hot ground-truth source label. The softmax formulation is important because it frames source detection as a competition among nodes, whereas a sigmoid output would treat each node's probability independently and is inappropriate for the single-source setting.

\subsection{GNNs for Temporal Source Detection}

Static GNNs applied to the aggregated graph $G_a$ are inherently limited by the loss of temporal ordering: two nodes may have identical degree and neighbourhood structure in $G_a$ while having very different causal reachability sets in the temporal network $\mathcal{G}$. The goal of temporal GNN architectures is to encode the time-varying contact structure in a way that is informative for source inference. Three broad strategies have emerged in the literature and are evaluated in this thesis. The first strategy, employed by the Backtracking Network of Ru et al.~\cite{Ru2023}, augments the edges of the aggregated graph with binary activation vectors---referred to as edge textures---that record at which discrete time steps each contact was active; this preserves the scalability of standard graph convolution while providing the model with access to temporal co-occurrence patterns (Section~\ref{ssst:kernelmodel}). The second strategy, used by the De Bruijn Graph Neural Network of Qarkaxhija et al.~\cite{Qarkaxhija2022}, lifts the temporal network into a higher-order state space whose nodes represent ordered sequences of contacts, thereby encoding causal walk structure as an explicit graph topology on which standard message passing can operate (Section~\ref{ssst:debruijnmodel}). The third strategy combines the Temporal Event Graph representation of Saram\"{a}ki et al.~\cite{Saramaki2019}---which maps the temporal network into a directed acyclic graph whose nodes are contact events and whose edges encode time-respecting sequential relationships---with the DAG convolutional network of Rey et al.~\cite{Rey2025}, enabling spectral learning directly on the acyclic causal structure (Section~\ref{ssst:dagmodel}). Each strategy encodes a different inductive bias about how temporal information should be represented and aggregated, and a primary contribution of this thesis is a systematic comparison of the detection accuracy achieved by each approach under identical experimental conditions.
